% ----------------------
\subsection{Section 111 (6 August 1836)}
% ----------------------
	\label{DC111}
	
	% -----
	\subsubsection{Historical Background}
	% -----
		\label{DC111-HB}
		
		% --
		\paragraph{JSP}
		% --
		
			``On 25 July 1836, after writing two letters concerning church members in Missouri, JS, Hyrum Smith, Oliver Cowdery, and Sidney Rigdon left Kirtland, Ohio, to travel to the eastern United States, briefly visiting New York City and Boston and staying in Salem, Massachusetts, before returning to Kirtland in mid-September. Salem, which was officially incorporated as a city in May 1836, was described by Oliver Cowdery as a `pleasantly situated town with fifteen thousand inhabitants.' Located on Massachusetts Bay about fifteen miles north of the larger city of Boston, Salem's busy port held a prominent place in the domestic and international commercial shipping trade of the United States. JS and his three companions arrived in Salem on 5 August and rented a house on Union Street for three weeks. The house may have been where JS dictated this revelation a day later.%
				\footnote{%
					% \label{}%
					HDDC 1468 says, ``While Joseph Smith was in Salem, he recieved a visit from Brigham Young and Lyman E. Johnson. While they were with him, Joseph received this revelation.'' The reference is then to \textit{Millennial Star} 25 July 1863, page 472.
					}

			``No known contemporary documents specify church leaders' reasons for visiting the eastern United States, and few records discuss the trip. The main contemporary sources of information are two letters written by Oliver Cowdery to his brother Warren A. Cowdery, which were published in the church's newspaper, and a letter JS wrote to his wife Emma while in Salem. Oliver Cowdery's letters indicate that their time in New England was spent preaching and occasionally visiting historic places. The four church leaders were likely motivated by a concern about Zion and the financial situation of the church, particularly a need to reduce debts of church leaders. The financial burden placed on them by finishing the House of the Lord in Kirtland and purchasing land in Ohio and Missouri had added significantly to the church's existing debts. Following a 2 April 1836 meeting at which JS and Cowdery were assigned to raise money to purchase land in Missouri, the men appear to have encountered difficulties in finding members willing to give their money or land for the cause of Zion. With the citizens of Clay County, Missouri, requesting that church members living there relocate, the need for temporal means to aid church members in Missouri grew even more pressing. JS and his colleagues may have raised money as they preached during their trip east in 1836, as it was not uncommon for missionaries to have the dual objectives of proselytizing and collecting funds for the church. This 6 August revelation addresses the church leaders' financial concerns. It informed the men that they would have the power to pay their debts and instructed them to `concern not yourselves about Zion' for there were people and money in Salem `for the benefit of Zion.'

			``Related to the revelation's statement that there was `much treasure' in Salem, two later accounts from individuals not directly involved in the journey state that JS traveled to the eastern United States to search for treasure or hidden money. In an 1843 pamphlet, sixteen-year-old dissident James C. Brewster briefly mentioned treasure hunting in relation to JS's 1836 trip. Ebenezer Robinson wrote an account in 1889, fifty-three years after JS's trip, that also linked the 1836 trip and searching for treasure---in fact, he claimed that the single objective of the trip was to look for hidden money in Salem. Robinson printed his account as the editor of the \textit{Return}, a publication for David Whitmer's Church of Christ. Robinson, who joined the Church of the Latter Day Saints in 1836 while working in the Kirtland printing office, stated in his reminiscences that Don Carlos Smith, who worked with him, told him that JS had learned about possible treasure from `a brother in the church, by the name of Burgess' who had come to Kirtland and `stated that a large amount of money had been secreted in the cellar of a certain house in Salem, Massachusetts, which had belonged to a widow, and he thought he was the only person now living, who had knowledge of it, or the location of the house.' Robinson claimed he was also told that Burgess met JS in Salem but that Burgess was unable to identify the house after so many years and left. Continuing their search, according to Robinson, JS and the three other men found and rented a house they thought contained the hidden money, but they were unsuccessful in finding it.

			``It is possible that JS had been told about hidden money in Salem and decided to pursue it to aid the church and relieve the financial and temporal pressure weighing down the branches in Kirtland and Missouri, and two contemporary documents may provide support for the statements of Brewster and Robinson. First, a promissory note was made out to a Jonathan Burgess in Salem, a tentative connection to the Burgess of Robinson's account. Second, JS mentioned looking for a specific house in Salem in his 19 August letter to Emma Smith. Robinson's account stated that JS rented the house and failed to find any treasure, but JS's letter to Emma reveals that he had not been able to rent or gain access to the house. While JS seemed hopeful the situation would change, the men left Salem only a few days later and offered no indication that they had rented or even visited the sought-after house, nor is there any evidence that they later returned.

			``Aside from alluding to `more treasures than one,' the revelation makes other references to the people of Salem and their significance to the growing church. In the three weeks following the revelation, JS and the others in the church presidency spent much of their time in Salem and in Boston preaching to the people. By the early 1840s, church leaders in Nauvoo focused on the proselytizing aspects of the 6 August 1836 revelation.%
				\footnote{%
					% \label{}%
					See HDDC 1468--1469 for more here, including a long quotation from Erastus Snow's journal on his missionary efforts related to Salem.
					}
			In 1841, Hyrum Smith and William Law of the First Presidency visited the eastern United States and left instructions at a church conference in Philadelphia for Erastus Snow and Benjamin Winchester to extend their missions and begin preaching in Salem. Snow recorded in his journal that Smith and Law `left with us a copy of a Revelation given about that people in 1836 which said the Lord had much people there whom he would gather into his kingdom in his Own due time and they thought the due time of the Lord had come.' Snow and Winchester arrived in Salem on 3 September 1841. After a week, Winchester returned to Philadelphia while Snow preached in Salem and the surrounding area. Snow organized the Salem branch on 5 March 1842, and by the end of his mission more than one hundred people had joined the church. When Snow and his family left in the fall of 1843, seventy-five members from `Boston and the eastern churches' traveled with them to Nauvoo.

			``The original text of the revelation has not been found, but four copies are extant. The version presented here comes from a diary kept by William W. Phelps. Phelps's diary also contains earlier JS revelations, Phelps's September 1835 patriarchal blessing, and later material from the Illinois and Utah eras of the church. Although Phelps was in Kirtland for the March 1836 dedication of the House of the Lord, he had returned to Missouri before the date of this revelation. He likely copied the revelation into his diary in the late 1830s or early 1840s. Based on textual comparison, Phelps's copy appears to be the earliest version, and it matches the text of later printed versions. Another version, made by Erastus Snow, was likely copied in 1841 from a copy that was left for him in Philadelphia by Hyrum Smith and William Law. Snow's copy matches the wording used in Phelps's inscription, with only minor exceptions. A third copy is in the Book of the Law of the Lord, inscribed by Robert B. Thompson in Nauvoo between 1840 and 1841. A fourth inscription of the revelation is found in volume B-1 of JS's history and was written by Willard Richards between 1842 and 1844. The version in JS's history is textually similar to both the Phelps and Book of the Law of the Lord inscriptions, with punctuation and spelling in the first six lines of the revelation that match the Book of the Law of the Lord inscription.

			``This revelation was not published in JS's lifetime. It first appeared in print in the 1850s with the printing of the `History of Joseph Smith' in the \textit{Deseret News} and the \textit{Millennial Star}. It was first included in the Doctrine and Covenants in the 1876 edition, and the canonized version follows the text of the Phelps inscription featured here. Significant differences between Phelps's copy and other early copies of this revelation are described in footnotes below.''
					
	% -----
	\subsubsection{Versions}
	% -----
		\label{DC111-V}
		
		See the \href{https://drive.google.com/file/d/1goAvNz8jS4R8slYfMG_db55Ty2z_vmgK/view?usp=sharing}{sections comparison} document.
	
		\begin{compactdesc}
			\item[JSP] \hyperref[EMS-1836-JS-6-AUG-1836]{6 August 1836}
			\item[BC] ---
			\item[DC 1835] ---
			\item[DC 1844] ---
			\item[TS] ---
			\item[DN] 3:3 (25 December 1852), 9
			\item[MS] 15:51 (17 December 1853), 822
		\end{compactdesc}

	% -----
	\subsubsection{Section 111:1} % *DC111-1 
	% -----
		\label{DC111-1}

		I, THE Lord your God, am not displeased with your coming this journey, notwithstanding your follies.

		% \begin{framed}
		% \scriptsize
			% \begin{compactdesc}
				% \item[JSP] 
				% % \item[BC] 
				% % \item[]
			% \end{compactdesc}
		% \end{framed}

	% -----
	\subsubsection{Section 111:2} % *DC111-2 
	% -----
		\label{DC111-2}

		I have much treasure in this city for you, for the benefit of Zion, and many people in this city, whom I will gather out in due time for the benefit of Zion, through your instrumentality.

		% \begin{framed}
		% \scriptsize
			% \begin{compactdesc}
				% \item[JSP] 
				% % \item[BC] 
				% % \item[]
			% \end{compactdesc}
		% \end{framed}

	% -----
	\subsubsection{Section 111:3} % *DC111-3 
	% -----
		\label{DC111-3}

		Therefore, it is expedient that you should form acquaintance with men in this city, as you shall be led, and as it shall be given you.

		% \begin{framed}
		% \scriptsize
			% \begin{compactdesc}
				% \item[JSP] 
				% % \item[BC] 
				% % \item[]
			% \end{compactdesc}
		% \end{framed}

	% -----
	\subsubsection{Section 111:4} % *DC111-4 
	% -----
		\label{DC111-4}

		And it shall come to pass in due time that I will give this city into your hands, that you shall have power over it, insomuch that they shall not discover your secret parts; and its wealth pertaining to gold and silver shall be yours.

		% \begin{framed}
		% \scriptsize
			% \begin{compactdesc}
				% \item[JSP] 
				% % \item[BC] 
				% % \item[]
			% \end{compactdesc}
		% \end{framed}

	% -----
	\subsubsection{Section 111:5} % *DC111-5 
	% -----
		\label{DC111-5}

		Concern not yourselves about your debts, for I will give you power to pay them.

		% \begin{framed}
		% \scriptsize
			% \begin{compactdesc}
				% \item[JSP] 
				% % \item[BC] 
				% % \item[]
			% \end{compactdesc}
		% \end{framed}

	% -----
	\subsubsection{Section 111:6} % *DC111-6 
	% -----
		\label{DC111-6}

		Concern not yourselves about Zion, for I will deal mercifully with her.

		% \begin{framed}
		% \scriptsize
			% \begin{compactdesc}
				% \item[JSP] 
				% % \item[BC] 
				% % \item[]
			% \end{compactdesc}
		% \end{framed}

	% -----
	\subsubsection{Section 111:7} % *DC111-7 
	% -----
		\label{DC111-7}

		Tarry in this place, and in the regions round about;

		% \begin{framed}
		% \scriptsize
			% \begin{compactdesc}
				% \item[JSP] 
				% % \item[BC] 
				% % \item[]
			% \end{compactdesc}
		% \end{framed}

	% -----
	\subsubsection{Section 111:8} % *DC111-8 
	% -----
		\label{DC111-8}

		And the place where it is my will that you should tarry, for the main, shall be signalized unto you by the peace and power of my Spirit, that shall flow unto you.

		% \begin{framed}
		% \scriptsize
			% \begin{compactdesc}
				% \item[JSP] 
				% % \item[BC] 
				% % \item[]
			% \end{compactdesc}
		% \end{framed}

	% -----
	\subsubsection{Section 111:9} % *DC111-9 
	% -----
		\label{DC111-9}

		This place you may obtain by hire.  And inquire diligently concerning the more ancient inhabitants and founders of this city;

		% \begin{framed}
		% \scriptsize
			% \begin{compactdesc}
				% \item[JSP] 
				% % \item[BC] 
				% % \item[]
			% \end{compactdesc}
		% \end{framed}

	% -----
	\subsubsection{Section 111:10} % *DC111-10 
	% -----
		\label{DC111-10}

		For there are more treasures than one for you in this city.

		% \begin{framed}
		% \scriptsize
			% \begin{compactdesc}
				% \item[JSP] 
				% % \item[BC] 
				% % \item[]
			% \end{compactdesc}
		% \end{framed}

	% -----
	\subsubsection{Section 111:11} % *DC111-11 
	% -----
		\label{DC111-11}

		Therefore, be ye as wise as serpents and yet without sin; and I will order all things for your good, as fast as ye are able to receive them.  Amen.

		% \begin{framed}
		% \scriptsize
			% \begin{compactdesc}
				% \item[JSP] 
				% % \item[BC] 
				% % \item[]
			% \end{compactdesc}
		% \end{framed}